\clearpage\nobalance
\section*{Author Notes}
\subsection*{Acknowledgements}
I thank Prof. Luyao Zhang, Program Chair, and Prof. Ken Rogerson, Program Discussant, of the \emph{2056 Nobel and Turing Laureate Program Launch Workshop}, held at Duke Kunshan University on September 7, 2026. I also thank [full name of teammate 1] and [full name of teammate 2], my teammates in [session number and full session title], and the rest of the COMSCI/ECON 206 Autumn 2026 Session 1 class for their questions and discussion. [Identify a specific contribution you actually used.] I am responsible for the proposal and remaining errors.
\subsection*{Individual contribution}
[State what you personally designed, implemented, wrote, checked, and adapted. Identify shared code or other contributions accurately. Program roles and feedback do not automatically imply coauthorship.]

\bibliographystyle{ACM-Reference-Format}\bibliography{references}
\appendix
\section{Technical details and reproducibility}
[Supply derivations, inputs, package versions, seed where relevant, actual outputs and limitations. Identify repository commit, notebook cells and demo version. Include the fresh-run record.]
\subsection{AI-use disclosure}
AI-Assisted drafting and debugging are permitted after initial independent reasoning. Initial reasoning and handwritten reflection are Human-Only. Record tool/model, date, purpose, significant prompts, accepted or rejected suggestions, human changes and verification. If none, state ``No generative AI was used.'' The human author is responsible for every claim, citation and computation.

\section{Cumulative Intellectual Development}
This proposal develops my Intellectual Statement II through the \emph{2056 Nobel and Turing Laureate Program Launch Workshop}, held on September 7, 2026 at Duke Kunshan University, IB 2050. Program Chair: Prof. Luyao Zhang. Program Discussant: Prof. Ken Rogerson. My session: [number and full title]. My role: [assigned role]. My teammates: [full name 1] and [full name 2]. I acknowledge their contributions and the rest of the class's questions and discussion.

[Trace an earlier claim, a specific dated comment or piece of evidence, and the resulting revision with its intellectual consequence. Do not write only a chronological diary. Retain and revise the 2056 contribution statement.]

Appendix~\ref{app:abstract} contains the structured abstract in the cumulative reflection format. Update its eight rows and preserve the prior version. The Expected 2056 Nobel Prize and Turing Award contribution statement in Section 5 must follow from the foundation, current-era boundary and interdisciplinary milestones.

\section{Field and open-source inspiration}
The \emph{Joint Interdisciplinary Field Trip---Technology, Computation, Culture \& Society in Action} took place on September 4, 2026. Record your actual observations at \textbf{Tencent Shanghai Office} and \textbf{Shanghai Science and Technology Museum}. State attendance accurately.

[Retain one industry photograph and one public-sector photograph from the earlier assignment. Each caption gives exact site, September 4, 2026 visit date, photographer/source, activity or exhibit, and relevance. Distinguish observation from interpretation; connect each setting to an open-source inspiration, subsequent research, and proposed practical value. If you lacked first-hand access, identify an attributed alternative instead of claiming attendance.]

\section{Review response and revision record}
[Retain v1, both peer reviews, author rebuttal, reviewer follow-up and v2. Use the three review themes: strengths to retain, constructive concerns to address, and questions to clarify. Explain your decisions, changed locations and actual verification results. Justify any retained approach or deferred work.]

\clearpage\onecolumn
\section{Structured abstract and cumulative proposal map}\label{app:abstract}
Update this structured abstract with each cumulative reflection. Replace each prompt with a brief, concrete entry and retain earlier versions. The first seven rows summarize the research; the final row records feedback and revision. Detailed optional prompts are in the separate Scaffolding Companion.
\par\medskip
\begin{table}[H]
\caption{Appendix E structured abstract template; retain the eight-row proposal structure.}
\label{tab:abstract-map}
\renewcommand{\arraystretch}{1.18}
\begin{tabularx}{\textwidth}{@{}p{0.27\textwidth}Y@{}}
\toprule
\textbf{Proposal element} & \textbf{Current statement / change} \\
\midrule
Background & State the strategic setting, established knowledge and unresolved issue. Cite the key foundation. \\
Motivation and application & Explain why the issue matters now, the application scenario, affected people and the need the project addresses. \\
Three connected questions & State one economics question, one computer-science question and one behavioral-science question. End with the integrated question that requires all three. \\
Method and model assumptions & Specify players, actions or strategies, incentives, timing and information; then the model, baseline, computation and evidence needed. \\
Results or expected results & State what was actually obtained and label its evidence type. Separately state expected results, the test and what finding would challenge the claim. \\
Intellectual merits & Explain what is learned beyond the cited foundation and why integrating the three disciplines changes the answer. \\
Practical impacts & Identify private- and public-sector value. Explain the open-access learning tool's intended contribution to SDG 4 and how its usefulness could be assessed. \\
Feedback received and revision made & Name the dated feedback source, earlier claim, actual revision and intellectual consequence. Preserve prior versions; state accurately if feedback is pending. \\
\bottomrule
\end{tabularx}
\end{table}
\noindent Keep field observations, photographs, literature notes, open-source inspirations and AI-use disclosure in the supporting sections. Distinguish observation, interpretation and proposed future evidence. The Open Science Statement groups GitHub and Colab; the Statement of Contribution to the UN's SDGs groups the Hugging Face learning tool.

